\documentclass[11pt,a4paper]{article}
\newcommand{\fichetype}{Corrigé détaillé}
% Thème A4 commun, librement modifiable.
\usepackage[T1]{fontenc}
\usepackage[utf8]{inputenc}
\usepackage[french]{babel}
\usepackage{lmodern,amsmath,amssymb,mathtools,tikz,pgfplots}
\usepackage[margin=2.2cm,headheight=15pt]{geometry}
\usepackage{xcolor,fancyhdr,enumitem,hyperref,needspace}
\usetikzlibrary{arrows.meta,positioning,calc}
\pgfplotsset{compat=1.18,width=11cm,height=6cm}
\definecolor{Ink}{HTML}{202124}
\definecolor{Blue}{HTML}{356B9E}
\definecolor{Muted}{HTML}{737579}
\color{Ink}
\renewcommand{\familydefault}{\sfdefault}
\setlength{\parindent}{0pt}
\setlength{\parskip}{6pt}
\setlist{itemsep=3pt,topsep=4pt}
\pagestyle{fancy}\fancyhf{}
\fancyhead[L]{\small\color{Muted}MATHS LYCÉE}
\fancyhead[R]{\small\color{Muted}2026--2027}
\fancyfoot[L]{\small\color{Muted}\fichetype}
\fancyfoot[R]{\small\thepage}
\renewcommand{\headrulewidth}{0.3pt}
\hypersetup{colorlinks=true,urlcolor=Blue,linkcolor=Blue}
\newcommand{\R}{\mathbb{R}}
\newcommand{\N}{\mathbb{N}}
\newcommand{\E}{\mathbb{E}}
\newcommand{\V}{\mathrm{V}}
\newcommand{\exercice}[3]{\Needspace{6\baselineskip}\vspace{0.5cm}{\color{Blue}\small\textbf{EXERCICE #1}\hfill #2}\par{\large\bfseries #3}\par}

\begin{document}
{\small\color{Blue}PREMIÈRE SPÉCIALITÉ}\par{\LARGE\bfseries La fonction exponentielle}\par{\large Corrigé détaillé}\par\vspace{0.2cm}\hrule\vspace{0.4cm}
\par\noindent\begin{minipage}{\linewidth}\exercice{1}{Application}{Simplifier}
{\small\color{Muted}\textbf{Énoncé.} Simplifier $e^{2x}e^{3-x}$ et $(e^{2x})^3/e^x$.}
\par\medskip
\textbf{Résolution.}\par
On additionne les exposants du produit : $e^{x+3}$. Dans la seconde expression, la puissance donne $e^{6x}$ puis le quotient donne $e^{5x}$.
\end{minipage}\par
\par\noindent\begin{minipage}{\linewidth}\exercice{2}{Application}{Équation}
{\small\color{Muted}\textbf{Énoncé.} Résoudre $e^{4x-2}=e^{x+7}$.}
\par\medskip
\textbf{Résolution.}\par
La stricte croissance donne $4x-2=x+7$, soit $3x=9$, donc $x=3$. Vérification : les deux exposants valent $10$.
\end{minipage}\par
\par\noindent\begin{minipage}{\linewidth}\exercice{3}{Application}{Inéquation}
{\small\color{Muted}\textbf{Énoncé.} Résoudre $e^{1-3x}>e^4$.}
\par\medskip
\textbf{Résolution.}\par
On compare les exposants : $1-3x>4$, donc $-3x>3$. Diviser par $-3$ inverse le sens : $x<-1$.
\end{minipage}\par
\par\noindent\begin{minipage}{\linewidth}\exercice{4}{Entraînement}{Dériver}
{\small\color{Muted}\textbf{Énoncé.} Dériver $3e^{2x-1}$ et $(x+1)e^{-x}$.}
\par\medskip
\textbf{Résolution.}\par
La première dérivée vaut $6e^{2x-1}$. Pour la seconde, le produit donne $e^{-x}-(x+1)e^{-x}=-xe^{-x}$. On a utilisé la dérivée de $e^{-x}$, qui vaut $-e^{-x}$.
\end{minipage}\par
\par\noindent\begin{minipage}{\linewidth}\exercice{5}{Entraînement}{Maximum}
{\small\color{Muted}\textbf{Énoncé.} Étudier $f(x)=x e^{-2x}$ sur $[0,3]$.}
\par\medskip
\textbf{Résolution.}\par
$f\prime=(1-2x)e^{-2x}$. Le second facteur est positif, le premier change de signe en $1/2$. La fonction croît puis décroît. Maximum $f(1/2)=1/(2e)$, et valeurs aux bornes $f(0)=0$, $f(3)=3e^{-6}$.
\end{minipage}\par
\par\noindent\begin{minipage}{\linewidth}\exercice{6}{Entraînement}{Population}
{\small\color{Muted}\textbf{Énoncé.} Pour $N(t)=500e^{0,04t}$, donner $N(0)$, $N\prime(t)$ et la raison de $N(n)$.}
\par\medskip
\textbf{Résolution.}\par
$N(0)=500$. La dérivée est $20e^{0,04t}=0,04N(t)$. Aux temps entiers, $N(n+1)/N(n)=e^{0,04}\approx1,04081$. L’augmentation par unité est donc d’environ $4,08\,\%$.
\end{minipage}\par
\par\noindent\begin{minipage}{\linewidth}\exercice{7}{Approfondissement}{Seuil numérique}
{\small\color{Muted}\textbf{Énoncé.} Décrire un algorithme cherchant le premier entier $n$ tel que $100e^{0,1n}>200$.}
\par\medskip
\textbf{Résolution.}\par
Initialiser $n=0$. Tant que $100e^{0,1n}\le200$, augmenter $n$ de un. On obtient $n=7$ : $100e^{0,6}\approx182,21$ et $100e^{0,7}\approx201,38$. La croissance garantit qu’il s’agit du premier entier.
\end{minipage}\par
\par\noindent\begin{minipage}{\linewidth}\exercice{8}{Approfondissement}{Une fausse propriété}
{\small\color{Muted}\textbf{Énoncé.} Un élève affirme $e^{a+b}=e^a+e^b$. Donner un contre-exemple et la propriété correcte.}
\par\medskip
\textbf{Résolution.}\par
Prendre $a=b=0$. Le membre de gauche vaut $1$, celui de droite $2$ : l’affirmation est fausse. La relation correcte est $e^{a+b}=e^ae^b$. Un contre-exemple suffit à réfuter un énoncé universel.
\end{minipage}\par
\end{document}
